\documentclass[a4paper,11pt]{book}
%\documentclass[a4paper,twoside,11pt,titlepage]{book}
\usepackage{listings}
\usepackage[utf8]{inputenc}
\usepackage[spanish]{babel}

% \usepackage[style=list, number=none]{glossary} %
%\usepackage{titlesec}
%\usepackage{pailatino}

%\decimalpoint
\usepackage{dcolumn}
\newcolumntype{.}{D{.}{\esperiod}{-1}}
\makeatletter
%\addto\shorthandsspanish{\let\esperiod\es@period@code}
\makeatother


%\usepackage[chapter]{algorithm}
\RequirePackage{verbatim}
%\RequirePackage[Glenn]{fncychap}
\usepackage{fancyhdr}
\usepackage{graphicx}
\usepackage{afterpage}

\usepackage{longtable}

\usepackage[pdfborder={000}]{hyperref} %referencia

% ********************************************************************
% Re-usable information
% ********************************************************************
\newcommand{\myTitle}{NO-INVENTORY\xspace}
\newcommand{\myDegree}{Grado en Ingeniería Informática\xspace}
\newcommand{\myName}{César Hugo Bárzano Cruz\xspace}
\newcommand{\myProf}{Nombre Apllido1 Apellido2 (tutor1)\xspace}
\newcommand{\myOtherProf}{Nombre Apllido1 Apellido2 (tutor2)\xspace}
%\newcommand{\mySupervisor}{Put name here\xspace}
\newcommand{\myFaculty}{Escuela Técnica Superior de Ingenierías Informática y de
Telecomunicación\xspace}
\newcommand{\myFacultyShort}{E.T.S. de Ingenierías Informática y de
Telecomunicación\xspace}
\newcommand{\myDepartment}{Departamento de Arquitectura y Tecnología de los Computadores\xspace}
\newcommand{\myUni}{\protect{Universidad de Granada}\xspace}
\newcommand{\myLocation}{Granada\xspace}
\newcommand{\myTime}{\today\xspace}
\newcommand{\myVersion}{Version 0.1\xspace}


\hypersetup{
pdfauthor = {\myName hugobarzano@correo.ugr.es},
pdftitle = {\myTitle},
pdfsubject = {},
pdfkeywords = {Gestión, Almacén , Inventario, Eficiencia, Ahorro, Informes,Web, Android },
pdfcreator = {LaTeX con el paquete TEXmaker},
pdfproducer = {pdflatex}
}

%\hyphenation{}


%\usepackage{doxygen/doxygen}
%\usepackage{pdfpages}
\usepackage{url}
\usepackage{colortbl,longtable}
\usepackage[stable]{footmisc}
%\usepackage{index}

%\makeindex
%\usepackage[style=long, cols=2,border=plain,toc=true,number=none]{glossary}
% \makeglossary

% Definición de comandos que me son tiles:
%\renewcommand{\indexname}{Índice alfabético}
%\renewcommand{\glossaryname}{Glosario}

\pagestyle{fancy}
\fancyhf{}
\fancyhead[LO]{\leftmark}
\fancyhead[RE]{\rightmark}
\fancyhead[RO,LE]{\textbf{\thepage}}
\renewcommand{\chaptermark}[1]{\markboth{\textbf{#1}}{}}
\renewcommand{\sectionmark}[1]{\markright{\textbf{\thesection. #1}}}

\setlength{\headheight}{1.5\headheight}

\newcommand{\HRule}{\rule{\linewidth}{0.5mm}}
%Definimos los tipos teorema, ejemplo y definición podremos usar estos tipos
%simplemente poniendo \begin{teorema} \end{teorema} ...
\newtheorem{teorema}{Teorema}[chapter]
\newtheorem{ejemplo}{Ejemplo}[chapter]
\newtheorem{definicion}{Definición}[chapter]

\definecolor{gray97}{gray}{.97}
\definecolor{gray75}{gray}{.75}
\definecolor{gray45}{gray}{.45}
\definecolor{gray30}{gray}{.94}

\lstset{ frame=Ltb,
     framerule=0.5pt,
     aboveskip=0.5cm,
     framextopmargin=3pt,
     framexbottommargin=3pt,
     framexleftmargin=0.1cm,
     framesep=0pt,
     rulesep=.4pt,
     backgroundcolor=\color{gray97},
     rulesepcolor=\color{black},
     %
     stringstyle=\ttfamily,
     showstringspaces = false,
     basicstyle=\scriptsize\ttfamily,
     commentstyle=\color{gray45},
     keywordstyle=\bfseries,
     %
     numbers=left,
     numbersep=6pt,
     numberstyle=\tiny,
     numberfirstline = false,
     breaklines=true,
   }

% minimizar fragmentado de listados
\lstnewenvironment{listing}[1][]
   {\lstset{#1}\pagebreak[0]}{\pagebreak[0]}

\lstdefinestyle{CodigoC}
   {
	basicstyle=\scriptsize,
	frame=single,
	language=C,
	numbers=left
   }
\lstdefinestyle{CodigoC++}
   {
	basicstyle=\small,
	frame=single,
	backgroundcolor=\color{gray30},
	language=C++,
	numbers=left
   }


\lstdefinestyle{Consola}
   {basicstyle=\scriptsize\bf\ttfamily,
    backgroundcolor=\color{gray30},
    frame=single,
    numbers=none
   }


\newcommand{\bigrule}{\titlerule[0.5mm]}


%Para conseguir que en las páginas en blanco no ponga cabecerass
\makeatletter
\def\clearpage{%
  \ifvmode
    \ifnum \@dbltopnum =\m@ne
      \ifdim \pagetotal <\topskip
        \hbox{}
      \fi
    \fi
  \fi
  \newpage
  \thispagestyle{empty}
  \write\m@ne{}
  \vbox{}
  \penalty -\@Mi
}
\makeatother

\usepackage{pdfpages}
\begin{document}
\begin{titlepage}
 
 
\newlength{\centeroffset}
\setlength{\centeroffset}{-0.5\oddsidemargin}
\addtolength{\centeroffset}{0.5\evensidemargin}
\thispagestyle{empty}

\noindent\hspace*{\centeroffset}\begin{minipage}{\textwidth}

\centering
\includegraphics[width=0.9\textwidth]{imagenes/logo_ugr.jpg}\\[1.4cm]

\textsc{ \Large TRABAJO FIN DE GRADO\\[0.2cm]}
\textsc{ INGENIERÍA EN INFORMÁTICA}\\[1cm]
% Upper part of the page
% 
% Title
{\Huge\bfseries NO-INVENTORY\\
}
\noindent\rule[-1ex]{\textwidth}{3pt}\\[3.5ex]
{\large\bfseries Sistema de gestión para inventarios}
\end{minipage}

\vspace{2.5cm}
\noindent\hspace*{\centeroffset}\begin{minipage}{\textwidth}
\centering

\textbf{Autor}\\ {César Hugo Bárzano Cruz}\\[2.5ex]
\textbf{Directores}\\
{Nombre Apellido1 Apellido2 (tutor1)\\
Nombre Apellido1 Apellido2 (tutor2)}\\[2cm]
\includegraphics[width=0.3\textwidth]{imagenes/etsiit_logo.png}\\[0.1cm]
\textsc{Escuela Técnica Superior de Ingenierías Informática y de Telecomunicación}\\
\textsc{---}\\
Granada, mes de 201
\end{minipage}
%\addtolength{\textwidth}{\centeroffset}
%\vspace{\stretch{2}}
\end{titlepage}



\chapter*{}
%\thispagestyle{empty}
%\cleardoublepage

%\thispagestyle{empty}

\begin{titlepage}
 
 
\setlength{\centeroffset}{-0.5\oddsidemargin}
\addtolength{\centeroffset}{0.5\evensidemargin}
\thispagestyle{empty}

\noindent\hspace*{\centeroffset}\begin{minipage}{\textwidth}

\centering
%\includegraphics[width=0.9\textwidth]{imagenes/logo_ugr.jpg}\\[1.4cm]

%\textsc{ \Large PROYECTO FIN DE CARRERA\\[0.2cm]}
%\textsc{ INGENIERÍA EN INFORMÁTICA}\\[1cm]
% Upper part of the page
% 

 \vspace{3.3cm}

%si el proyecto tiene logo poner aquí
\includegraphics{imagenes/logo.png} 
 \vspace{0.5cm}

% Title

{\Huge\bfseries Título del proyecto\\
}
\noindent\rule[-1ex]{\textwidth}{3pt}\\[3.5ex]
{\large\bfseries Subtítulo del proyecto.\\[4cm]}
\end{minipage}

\vspace{2.5cm}
\noindent\hspace*{\centeroffset}\begin{minipage}{\textwidth}
\centering

\textbf{Autor}\\ {Nombre Apellido1 Apellido2 (alumno)}\\[2.5ex]
\textbf{Directores}\\
{Nombre Apellido1 Apellido2 (tutor1)\\
Nombre Apellido1 Apellido2 (tutor2)}\\[2cm]
%\includegraphics[width=0.15\textwidth]{imagenes/tstc.png}\\[0.1cm]
%\textsc{Departamento de Teoría de la Señal, Telemática y Comunicaciones}\\
%\textsc{---}\\
%Granada, mes de 201
\end{minipage}
%\addtolength{\textwidth}{\centeroffset}
\vspace{\stretch{2}}

 
\end{titlepage}






\cleardoublepage
\thispagestyle{empty}

\begin{center}
{\large\bfseries NO-INVENTORY: Sistema de gestión para inventarios}\\
\end{center}
\begin{center}
César Hugo Bárzano Cruz\\
\end{center}

%\vspace{0.7cm}
Gestión, Almacén , Activos, Eficiencia, Ahorro, Informes, 
\noindent{\textbf{Palabras clave}: Gestión\_clave1, Almacen\_clave2, Activos\_clave3, Eficiencia\_clave4, Ahorro\_clave5 Informes \_clave6}\\

\vspace{0.7cm}
\noindent{\textbf{Resumen}}\\

Este proyecto surge a raíz de un problema real presentado por la Oficina de Software Libre. Desde principios de 2012 y a partir de un acuerdo con la Unidad de Calidad, la oficina se encarga de recoger material informático procedente de los distintivos organismos de la Universidad de Granada. Dicho material ha alcanzado una gran cantidad, por lo que es necesario un sistema para gestionarlo. 
Partiendo de este problema, existen necesidades similares en empresas, organizaciones, instituciones, comercios y todo tipo de negocio con un almacén de activos.

En dichos almacenes, las tareas de gestión suelen ser ineficientes, complejas y costosas.  El tiempo y dinero que conllevan estas tareas suele ser un factor a tener en cuenta, ya que una gran mayoría de empresas, con almacenes pequeños o medianos, no utilizan un sistema comercial debido a los costes que suponen su mantenimiento. Este problema, sigue presente en organizaciones con almacenes de gran tamaño, ya que  el número de empleados y el tiempo necesario para las tareas básicas de catalogación y administración suponen un coste a tener en cuenta.

Otro problema con el que se encuentran los empleados de estas entidades es la generación de informes. Recopilar información de los activos del almacén con el objetivo de representarla de manera adecuada a las necesidades de cada negocio, puede llegar a ser una tarea complicada. 
Como solución a estos problemas, surge el sistema de gestión NO-INVENTORY que pretende mejorar estas tareas, con el objetivo de ahorrar tiempo y dinero al cliente, automatizando y facilitando las tareas de gestión. 
La piedra angular del sistema es una plataforma web alojada en la nube por lo que no es necesaria la instalación de ningún software adicional en las máquinas del cliente que deseé comenzar a utilizarlo, solo es necesario un navegador.  En dicha plataforma, las tareas de administración se realizan de manera intuitiva permitiendo a los distintos empleados trabajar de forma cooperativa dentro del entorno colaborativo que representa al almacén de su empresa, utilizando un sistema flexible y personalizable en función de las características con las que se quieran clasificar los elementos. Facilita la agrupación de elementos con o sin propiedades comunes en colecciones denominadas catálogos en función de las necesidades de cada cliente y utiliza dichas colecciones para la generación automática de informes y gráficos representativos del estado del almacén. En función de la cantidad de datos o de la integridad de los mismo, la infraestructura subyacente del proyecto, permitiría correr instancias aisladas del sistema, para dar servicio solo ha esa empresa, aprovechando así los recursos de la plataforma al máximo. 

El sistema cuenta con una aplicación Android como extensión para realizar tareas de catalogación y clasificación dentro del propio almacén. La ventaja de esto reside en que hoy en día casi todo el mundo cuenta con un smartphone, dando la posibilidad a los empleados de llevar la gestión del almacén con sus propios dispositivos.

En función de los activos que formen el almacén, y del presupuesto que se quiera dedicar a etiquetar y clasificar cada objeto, el cliente puede decidir que método utilizar, ya que la funcionalidad de la aplicación Android es la de leer y escribir los identificadores de cada elemento, con soporte para:
\begin{enumerate}
\item Códigos de Barras
\item Códigos QR
\item Etiquetas NFC 
\end{enumerate}

En resumen, el sistema pretende reducir tiempo y dinero a las organizaciones que decidan utilizarlo para la gestión de sus almacenes con el objetivo  de optimizar las tareas de los empleados consiguiendo así un mayor rendimiento. 
Por último, resaltar que todo lo relativo al desarrollo del proyecto, tanto código como documentación, está liberado en el sistema de control de versiones Github, bajo una licencia GPL3 para que pueda ser utilizado o mejorado por la comunidad de software libre.


\cleardoublepage


\thispagestyle{empty}


\begin{center}
{\large\bfseries Project Title: Project Subtitle}\\
\end{center}
\begin{center}
First name, Family name (student)\\
\end{center}

%\vspace{0.7cm}
\noindent{\textbf{Keywords}: Keyword1, Keyword2, Keyword3, ....}\\

\vspace{0.7cm}
\noindent{\textbf{Abstract}}\\

Write here the abstract in English.

\chapter*{}
\thispagestyle{empty}

\noindent\rule[-1ex]{\textwidth}{2pt}\\[4.5ex]

Yo, \textbf{César Hugo Bárzano Cruz}, alumno de la titulación TITULACIÓN de la \textbf{Escuela Técnica Superior
de Ingenierías Informática y de Telecomunicación de la Universidad de Granada}, con DNI 77138361h, autorizo la
ubicación de la siguiente copia de mi Trabajo Fin de Grado en la biblioteca del centro para que pueda ser
consultada por las personas que lo deseen.

\vspace{6cm}

\noindent Fdo: Nombre Apellido1 Apellido2

\vspace{2cm}

\begin{flushright}
Granada a X de mes de 201 .
\end{flushright}


\chapter*{}
\thispagestyle{empty}

\noindent\rule[-1ex]{\textwidth}{2pt}\\[4.5ex]

D. \textbf{Nombre Apellido1 Apellido2 (tutor1)}, Profesor del Área de XXXX del Departamento YYYY de la Universidad de Granada.

\vspace{0.5cm}

D. \textbf{Nombre Apellido1 Apellido2 (tutor2)}, Profesor del Área de XXXX del Departamento YYYY de la Universidad de Granada.


\vspace{0.5cm}

\textbf{Informan:}

\vspace{0.5cm}

Que el presente trabajo, titulado \textit{\textbf{Título del proyecto, Subtítulo del proyecto}},
ha sido realizado bajo su supervisión por \textbf{Nombre Apellido1 Apellido2 (alumno)}, y autorizamos la defensa de dicho trabajo ante el tribunal
que corresponda.

\vspace{0.5cm}

Y para que conste, expiden y firman el presente informe en Granada a X de mes de 201 .

\vspace{1cm}

\textbf{Los directores:}

\vspace{5cm}

\noindent \textbf{Nombre Apellido1 Apellido2 (tutor1) \ \ \ \ \ Nombre Apellido1 Apellido2 (tutor2)}

\chapter*{Agradecimientos}
\thispagestyle{empty}

       \vspace{1cm}


Poner aquí agradecimientos...


%\frontmatter
\tableofcontents
%\listoffigures
%\listoftables

%
%\mainmatter
%\setlength{\parskip}{5pt}

%\input{capitulos/01_Introduccion}


\chapter{Introducción}


Este proyecto tiene el objetivo de solucionar un problema real. Desde principios de 2012, la Oficina de Software Libre (OSL) tiene un acuerdo con Unidad de Calidad de la Universidad de Granada. Dicho acuerdo establece que la oficina es la encargada de recoger el material  informático procedente de los distintivos organismos de la Universidad. Debido a la acumulación de material a lo largo de estos 5 años, la oficina de software libre necesita un sistema de gestión que le permita catalogar y clasificar este material.

El problema principal radica en el desconocimiento del estado en el que este material es recogido, es decir, la OSL no sabe si el material recibido funciona, no funciona, necesita ser actualizado o directamente es inservible y ha de ser reciclado. Hay que tener en cuenta el estado actual del almacén. Tras acumular material a lo largo de estos años, el almacén sufre una situación caótica debido al gran numero de equipos sin catalogar ni reflejar el estado en el que se encuentran. 

Por otra parte, el acuerdo mencionado anteriormente, obliga a la oficina a realizar informes cada cierto periodo de tiempo, reflejando el estado de las campañas de recogida, el estado del material recolectado, el peso...etc

El problema presentado no es algo aislado a la OSL, gran parte de las pequeñas y medianas empresas con almacenes que gestionar presentan problemas similares. Para las  empresas, no usar plataformas de gestión para sus productos supone carencia de estrategias de crecimiento, una inadecuada utilización de las tecnologías y conocimientos, pérdidas de recursos, debilidad financiera y deficiencias en toda la organización. 

En los últimos años, el análisis de la información asociada a los recursos o servicios de las empresas, ha tomado una gran importancia. El análisis de la información se ha convertido en un factor clave para la toma de decisiones en diversos sectores. El valor de la información ha alcanzado el nivel de otros recursos como por ejemplo los recursos financieros, materiales y humanos. El conocimiento del entorno y la información que podemos aprovechar de él son vitales para la toma de decisiones, ya sea fuera del entorno empresarial o dentro de él. Decisiones como por ejemplo entrar en nuevos mercados, proteger la empresa de agentes externos, innovación, desarrollo, inversiones...etc. Por ello las Tecnologías de Información forman un factor decisivo para dar lugar al crecimiento tanto de las PyMEs como de cualquier empresa.

Lo principales motivos por los que estas empresas no utilizan sistemas de gestión para sus productos se deben a los costes elevados, a la carencia de recursos, la falta de acceso a la información, etc... 

Las Tecnologías de la Información y Comunicaciones (TIC), se han convertido en la herramienta perfecta para llevar a cabo las tareas necesarias para aumentar 
la eficacia y eficiencia en el manejo de los recursos de la empresa. 


\chapter{Análisis y Estado del Arte}

\begin{figure}[H] 
\centering 
\includegraphics[scale=1.95]{imagenes/tstc.png}
\caption{ man apt-cache  }  
\end{figure}

Para realizar un mejor Analisis del sistema, vamos a desglosarlo en dos partes. Primero analizaremos la plataforma web y despues analizaremos la aplicación Android 

\section{Requisitos de Información Plataforma web}

Los requisitos de información se caracterizan por  reunir la información 	relevante para el cliente que debe gestionar y almacenar el sistema software.\par 

\textbf{RI-1. Item:} Representación de cada uno de los elementos en el sistema. 
Contenido: nombre, descripción, fecha de alta, ID, Localizador, organización de la que forma parte, el usuario que lo dio de alta, atributos. Los atributos son una tupla de 3 valores asignados en función de las preferencias con las que se quiera clasificar  cada elemento. Por ejemplo lugar del que procede, tipo de elemento y estado del elemento.


\textbf{RI-2. Catalogo:} Representación de cada una de las colecciones de elementos del sistema. Contenido: nombre, descripción, fecha de alta, ID, mensaje de aviso, fecha de expiración, tipo de catalogo,  organización de la que forma parte, el usuario que lo dio de alta, colección con las ID de los elementos que forman el catalogo.

\textbf{RI-3. Gráfico:} Información del estado del sistema en forma de grafico, en función de los catálogos creados por los distintos usuarios de una misma organización. Gráficos representativos basados en los atributos de los items que forman cada catalogo. Comparativa entre unidades y peso de esas unidades en el catalogo. 
Contenido: catalogo.

\textbf{RI-4. Informe:} Información del estado del sistema en forma de documento. Generado en función de los catálogos creados por los distintos usuarios de una misma organización y de los gráficos representativos. 
Contenido: nombre, fecha de creación, usuario, organización, datos del informe. 

\textbf{RI-5. Perfil de Usuario:} Información sensible del usuario que va a trabajar con el sistema. Contenido: nombre de usuario, contraseña, correo electrónico, organización. 

\textbf{RI-6. Preferencias:} Información de los atributos que se desean utilizar para clasificar los elementos del sistema dentro del entorno colaborativo de la organización.
Contenido: archivo csv de preferencias 1, archivo csv de preferencias 2, archivo csv de preferencias 3. 


\section{Requisitos Funcionales Plataforma web}

Como se define en la ingeniería de requisitos, los requisitos funcionales establecen los comportamientos del sistema.

\textbf{RF-1. Gestión de items:} El sistema deberá gestionar el ciclo de vida de items o elementos existentes en el sistema.
   

	RF-1.1. El sistema permitirá añadir items

	RF-1.2. El sistema permitirá buscar items

	RF-1.3. El sistema permitirá modificar items

	RF-1.4. El sistema permitirá eliminar items

	RF-1.5. El sistema permitirá agrupar items en colecciones


\textbf{RF-2. Gestión de catálogos:} El sistema deberá gestionar el ciclo de vida de conjuntos de items, estas colecciones reciben el nombre de catálogos o agrupaciones de items con atributos u objetivos similares.   

	RF-2.1. El sistema permitirá añadir catálogos

	RF-2.2. El sistema permitirá buscar catálogos

	RF-2.3. El sistema permitirá modificar catálogos

	RF-2.4. El sistema permitirá eliminar catálogos

	RF-2.5. El sistema permitirá limpiar catálogos

	RF-2.6. El sistema permitirá eliminar un item del catalogo

	RF-2.7. El sistema permitirá exportar los códigos qr de sus elementos

	RF-2.8. El sistema permitirá exportar los códigos de barras de sus elementos

	RF-2.9. El sistema permitirá controlar el peso de los elementos del catálogo

	RF-2.10. El sistema permitirá controlar el número de unidades del catálogo. 

\textbf{RF-3. Gestión de Gráficos:} El sistema deberá generar gráficos a raíz de los elementos que forman parte de un catálogo, mostrando una comparativa entre el numero de unidades y peso de dichas unidades. Los gráficos representativos se dividen en función de las preferencias del usuario. 

	RF-3.1. El sistema permitirá generar gráficos
	RF-3.1. El sistema permitirá exportar gráficos
	RF-3.1. El sistema permitirá seleccionar el tipo de gráfico	

	 
\textbf{RF-4. Gestión de Informes:} El sistema deberá generar informes significativos a raíz de los elementos que forman parte de un catálogo. El sistema proveerá al usuario de una plantilla editable para personalizar los aspectos e información necesaria para el informe. El sistema permitirá combinar distintos catálogos para generar un mismo informe.  El sistema permitirá incluir en los informes, los gráficos generados con anterioridad. 

	

	RF-4.1. El sistema permitirá generar informes
	RF-4.2. El sistema permitirá almacenar informes
	RF-4.3. El sistema permitirá modificar informes
	RF-4.4. El sistema permitirá eliminar informes
	RF-4.5. El sistema permitirá añadir catálogos a los informes
	RF-4.6. El sistema permitirá añadir gráficos a los informes
	RF-4.7. El sistema permitirá exportar informes

\textbf{RF-5. Gestión de Preferencias:} El sistema permitirá personalizar las preferencias de cada organización. Dichas preferencias son los atributos con los que van a catalogarse los items. Dichos atributos son los que se usarán para dotar de contenido a los catálogos. Estos atributos también son empleados para la generación de gráficos. 

	RF-5.1. El sistema permitirá cargar las preferencias desde fichero CSV
	RF-5.1. El sistema permitirá cargar preferencias por defecto 

\textbf{RF-6. Gestión de Usuarios:} El sistema deberá presentar un entorno colaborativo, clasificando a los usuarios por organizaciones de forma que distintos usuarios de una misma organización tengan acceso al mismo conjunto de recursos que define dicha organización. 

	RF-6.1. El sistema permitirá al administrador dar de alta a nuevos usuarios.
	RF-6.1. El sistema permitirá al administrador asignar perfiles de usuario.
	RF-6.2. El sistema permitirá iniciar sesión a los usuario.
	RF-6.2. El sistema permitirá cerrar sesión a los usuario.
	RF-6.2. El sistema registrará internamente las acciones realizadas por cada 	usuario, 	de cada organización. 

\section{Requisitos No Funcionales Plataforma web}
Los requisitos no funcionales, se refieren a todos los requisitos que no describen información a guardar, ni funciones a realizar por el sistema, sino características de funcionamiento.

\textbf{RNF-1} Necesitaremos que toda la información que se almacena sobre el sistema se mantenga segura, realizando copias de seguridad periódicas.

\textbf{RNF-2} Necesitaremos analizar la cantidad de elementos que una organización va a administrar. Si es considerada una cantidad importante, dicha organización recibirá servicio desde una imagen del sistema asilada (Docker). De esta forma conseguimos que el flujo de datos no afectes a la experiencia de los clientes.

\textbf{RNF-3} Necesitaremos disponer de una buena conexión a internet ya que el sistema da servicio desde la nube. 

DIAGRAMONESSSSSSSSSSSSSSSSSSSSSSSSSS


DIAGRAMAAAAAAAAAAAAAOOOOOOOOOOOOOOOOOOOOOOSSS

\section{Requisitos Funcionales Aplicación Android}
Los requisitos funcionales de la aplicación Android establecen los comportamientos o funcionalidades de la aplicación.

RF-1. Administración de items 
La aplicación deberá administrar los items que formen parte del almacen.   

	RF-1.1. La aplicación permitirá detectar items mediante texto o localizador
	RF-1.2. La aplicación permitirá detectar items mediante codigo QR
	RF-1.3. La aplicación permitirá detectar items mediante codigo de barras
	RF-1.4. La aplicación permitirá detectar items mediante etiquetas NFC
	RF-1.5. La aplicación permitirá escribir el localizador de un item en una etiqueta NFC
	RF-1.6. La aplicación permitirá consultar un item en detalle
	RF-1.7. La aplicación permitirá añadir un item detectado a un catalogo determinado
	RF-1.8. La aplicación permitirá eliminar items


RF-2. Administración de Catálogos 
La aplicación deberá administrar los catálogos que se utilicen para clasificar los items del almacén.   

	RF-2.1. La aplicación permitirá detectar catálogos mediante texto o ID
	RF-2.2. La aplicación permitirá añadir items  al catalogo mediante código QR
	RF-2.3. La aplicación permitirá añadir items al catalogo mediante código de barras
	RF-2.4. La aplicación permitirá añadir items al catalogo mediante etiquetas NFC
	RF-2.5. La aplicación permitirá eliminar items de un catalogo
	RF-2.6. La aplicación permitirá limpiar la lista de items de un catalogo
	
RF-3. Sesión de Usuario 
La aplicación deberá identificar a cada usuario de manera única para certificar las tareas de administración que se lleven acabo desde la aplicación. 

	RF-3.1. La aplicación permitirá al usuario iniciar sesión.
	RF-3.2. La aplicación permitirá al usuario cerrar sesión.


\chapter{Resolución del trabajo}

Como método de ingeniería del software decir que vamos a seguir la técnicas de modelo de prototipos rápido o también llamado modelado de prototipado rápido.

Ver \url{http://www.ecured.cu/index.php/Modelo_de_Prototipos}

\section{Recursos}

Decir los recursos humanos (autor y directores), hardware y software que se van a utilizar.

\section{Especificación de requisitos}

Decir que se partió de una especificación inicial de requisitos que a medida que se fueron implementando los prototipos se fue refinando posteriormente. Se puede poner la inicial y la final o solo la final indicando que se están poniendo los requisitos que finalmente tiene que tener el sistema.

Los requisitos se pueden referir a las necesidades del usuario del sistema (requisitos del usuario), a lo que tiene que hacer la aplicación (requisito funcional) o a cómo tiene que hacerlo (requisito no funcional). Ejemplo:

En este sistema un robot tiene que coger con sus pinzas un envase de medicamento y llevárselo a una persona anciana que por sí misma no puede recordar su medicación.

Requisitos del usuario:

RU1. La persona puede moverse libremente por una habitación donde está el robot.

RU2. La persona es capaz de coger el medicamento cuando se lo ofrece el robot.

RU3. La persona es capaz de tomarse el medicamento por sí misma.


Requisitos funcionales.

RF1. El robot mediante la cámara kinect debe poder localizar a la persona.

RF2. El robot conoce la posición de la mesa pues tiene un mapa de la habitación.

RF3. El robot debe identificar el medicamento correcto según un plan de medicación previamente establecido.

RF4. El robot debe poder coger el medicamento con sus pinzas.

etc...

Requisitos no funcionales

RNF1. El robot no puede atropellar ni dañar a la persona en ningún momento.

RNF2. La aplicación debe ejecutarse en entornos linux

RNF3. La aplicación debe utilizar pocos recursos para reaccionar con rapidez.

algo de la interfaz, como tratar posibles fallos, etc.


\section{Planificación}

Poner una tabla de tiempos con las planificación del proyecto diciendo cuando se tiene previsto alcanzar cada subobjetivo planteado. Con su correspondiente división en fases y tareas, y la posterior
comparación con los datos reales obtenidos tras realizar el proyecto. Entre las fases está la realización de los diferentes prototipos I, II y III por ejemplo.

Poner presupuesto según horas de trabajo estimadas.


\section{Análisis funcional}

A partir de aquí nos referimos solamente al prototipo final que da lugar a la aplicación final.

Hay que describir la funcionalidad que debe poseer el sistema para poder cumplir con los objetivos y requisitos que se han dicho previamente.La descripción de esta funcionalidad puede hacerse analizando las tareas (que aparecerán en la planificación) y estudiando la inter-relación entre ellas y sus conexiones.

Para la realización de este análisis se pueden utilizar Diagramas de Flujo para poder conocer generalmente un único punto de inicio y un único punto de término o en varios. 

\url{https://es.wikipedia.org/wiki/Diagrama_de_flujo}

 Se pueden plantear también casos de uso. Los diagramas de casos de uso sirven para describir la inter-relación entre el sistema y el usuario del mismo. Se pueden utilizar para plantear diferentes casos de interacción entre el robot y la persona y cómo tiene que reaccionar el sistema en cada caso.

\url{https://es.wikipedia.org/wiki/Diagrama_de_casos_de_uso}


\section{Implementación y pruebas}

Decir qué lenguaje de programación se ha utilizado y las tecnologías implicadas aunque se hayan comentado en el apartado de recursos. Justificar su uso (rendimiento, disponibilidad, etc.). Si se ha usado open source decirlo y explicar las ventajas.

Las pruebas se realizan para comprobar la verificación y validación del producto software. La verificación consiste es comprobar que el producto realiza lo que está programado, es decir la programación no tiene errores y funciona en todos los casos cumpliendo los requisitos. La validación tiene que ver con que cumpla con lo que espera el usuario.

Verificación y Validación: Conjunto de procesos de comprobación y
análisis que aseguran que el software que se desarrolla está acorde a su
especificación y cumple las necesidades de los clientes.

Verificación:
¿Estamos construyendo el producto correctamente?
e su especificacin.

Validación:
¿Estamos construyendo el producto correcto?
Comprueba que el software cumple las expectativas que el cliente espera
Importante: Nunca se va a poder demostrar que el software está
completamente libre de defectos

las pruebas que pueden utilizarse son muy diversas. Aconsejo centrarnos en pruebas de caja blanca y de caja negra.

Las pruebas de la caja blanca se centran en la estructura interna del programa para elegir los casos de prueba. El objetivo de estas pruebas consiste en probar todos los posibles casos de ejecución de la aplicación para comprobar que los datos se comportan de manera correcta internamente.

Decir que se han hecho las pruebas de caja blanca.

Las pruebas de caja negra son aquellas que se centran en las salidas y entradas de los módulos, sin atender a su comportamiento interno (comprobando mediante las pruebas de caja blanca). Las pruebas de caja negra garantizan la interconectividad entre los diferentes módulos de la aplicación, así como su correcto funcionamiento final.

Poner algunos casos de prueba de caja negra.

\chapter{Conclusiones y trabajo futuro}

Decir lo que se ha conseguido realizar comentando sus puntos fuertes y débiles.
Decir si se han alcanzado los objetivos específicos y el general propuesto y en qué grado.

Indicar las asignaturas del grado más relacionadas con la ejecución del TFG y cómo el TFG ha ayudado a afianzar los conocimientos adquiridos en el Grado.

Valoración personal si se quiere.


Avanzar algunas líneas de trabajo futuro para solucionar las debilidades detectadas o para conseguir nuevas funcionalidades interesantes.

\chapter{Bibliografía}

Poner las citas bibliográficas, direcciones de internet, etc.


\chapter{Anexo}

Al final de la memoria hay que añadir un anexo de una página o dos, explicando como se usa el software a modo de manual de usuario. Es decir, como se llaman los comandos, qué parámetros hay que darle, como se llaman los ficheros de datos de entrada, etc.



%
%\input{capitulos/02_EspecificacionRequisitos}
%
%\input{capitulos/03_Planificacion}
%
%\input{capitulos/04_Analisis}
%
%\input{capitulos/05_Diseno}
%
%\input{capitulos/06_Implementacion}
%
%\input{capitulos/07_Pruebas}
%
%\input{capitulos/08_Conclusiones}
%
%%\chapter{Conclusiones y Trabajos Futuros}
%
%
%%\nocite{*}
%\bibliography{bibliografia/bibliografia}\addcontentsline{toc}{chapter}{Bibliografía}
%\bibliographystyle{miunsrturl}
%
%\appendix
%\input{apendices/manual_usuario/manual_usuario}
%%\input{apendices/paper/paper}
%\input{glosario/entradas_glosario}
% \addcontentsline{toc}{chapter}{Glosario}
% \printglossary
\chapter*{}
\thispagestyle{empty}

\end{document}
